\chapter{Peubah Acak Sinambung}\label{ch:g12:contdist}

Waktu tunggu, pengukuran fisika, proporsi: banyak besaran acak
bernilai pada suatu kontinum, dan tak ada satu nilai pun yang berpeluang positif.
Peluangnya lalu dihitung dengan mengintegralkan sebuah \emph{\omterm{def:g12:contdist:density}{kerapatan}}. Bab
ini mengkaji \omterm{def:g12:contdist:uniform}{distribusi seragam}, eksponen dan normal, lalu memakai
\omterm{def:g12:contdist:normal}{distribusi normalnya} untuk mengukur ayunan jajak pendapat dan contoh.

\section{Kerapatan peluang}

\begin{definition}[Kerapatan, peubah acak sinambung]\label{def:g12:contdist:density}
Suatu \emph{kerapatan peluang}\index{kerapatan} pada selang $I$ adalah \omterm{def:g11:func:function}{fungsi}
$f$ yang \omterm{def:g12:limcont:continuity}{kontinu} dan tak negatif pada $I$ dengan $\int_I f(t)\,\dd t = 1$
(dan \omterm{def:g12:integ:area}{integral} atas $I$ yang tak terbatas dipahami sebagai limit
\omterm{def:g12:integ:area}{integral} atas selang terbatas yang membesar). \omterm{def:g12:randvar:rv}{Peubah acak} $X$ berkerapatan
$f$ jika untuk semua $a \leq b$ di $I$:
\[
\P(a \leq X \leq b) = \int_a^b f(t)\,\dd t .
\]
\end{definition}

Peluangnya adalah \emph{luas di bawah kurva \omterm{def:g12:contdist:density}{kerapatannya}}. Khususnya
$\P(X = a) = \int_a^a f = 0$ untuk setiap nilai tunggal $a$: jadi hanya selang
yang memikul peluang, dan $\P(a \leq X \leq b) = \P(a < X < b)$.

\begin{definition}[Nilai harapan dan ragam]\label{def:g12:contdist:exp}
Untuk $X$ yang berkerapatan $f$ pada $I$:
\[
\E(X) = \int_I t\,f(t)\,\dd t, \qquad
\V(X) = \int_I \bigl(t - \E(X)\bigr)^2 f(t)\,\dd t
= \E(X^2) - \E(X)^2 .
\]
\end{definition}

Kaidah pada Bab 14--15 (kelinearan harapan, keaditifan
\omterm{def:g12:randvar:exp}{ragam} bagi peubah yang \omterm{def:g12:sums:indep}{saling bebas}, Bienaymé--Chebyshev) tetap berlaku;
buktinya, dengan \omterm{def:g12:integ:area}{integral} menggantikan jumlah, diterima saja pada tingkat ini.

\section{Distribusi seragam}

\begin{definition}[Distribusi seragam]\label{def:g12:contdist:uniform}
Peubah $X$ mengikuti \emph{distribusi seragam}\index{distribusi seragam}
$\mathcal U\intcc ab$ jika \omterm{def:g12:contdist:density}{kerapatannya} tetap,
$f(t) = \frac{1}{b-a}$ pada $\intcc{a}{b}$. Maka untuk
$\intcc{c}{d} \subseteq \intcc ab$,
$\P(c \leq X \leq d) = \frac{d - c}{b - a}$: yaitu peluang yang sebanding dengan
panjangnya.
\end{definition}

\begin{proposition}\label{prop:g12:contdist:uniformmoments}
Jika $X \sim \mathcal U\intcc ab$, maka
$\E(X) = \dfrac{a+b}{2}$ dan $\V(X) = \dfrac{(b-a)^2}{12}$.
\end{proposition}

\begin{proof}
$\E(X) = \int_a^b \frac{t}{b-a}\dd t
= \frac{b^2 - a^2}{2(b-a)} = \frac{a+b}{2}$. Untuk \omterm{def:g12:randvar:exp}{ragamnya},
$\E(X^2) = \frac{b^3 - a^3}{3(b-a)} = \frac{a^2 + ab + b^2}{3}$, dan
\[
\V(X) = \frac{a^2 + ab + b^2}{3} - \frac{(a+b)^2}{4}
= \frac{4a^2 + 4ab + 4b^2 - 3a^2 - 6ab - 3b^2}{12}
= \frac{(b - a)^2}{12}. \qedhere
\]
\end{proof}

\section{Distribusi eksponen}

\begin{definition}[Distribusi eksponen]\label{def:g12:contdist:expo}
Untuk $\lambda > 0$, peubah $X$ mengikuti \emph{distribusi
eksponen}\index{distribusi eksponen} $\mathcal E(\lambda)$ jika
\omterm{def:g12:contdist:density}{kerapatannya} pada $\intco{0}{+\infty}$ adalah
\[
f(t) = \lambda\,\eu^{-\lambda t}.
\]
Maka $\P(X \leq x) = 1 - \eu^{-\lambda x}$ dan
$\P(X > x) = \eu^{-\lambda x}$ untuk $x \geq 0$.
\end{definition}

\begin{proposition}\label{prop:g12:contdist:expomoments}
Jika $X \sim \mathcal E(\lambda)$:
$\E(X) = \dfrac{1}{\lambda}$ dan $\V(X) = \dfrac{1}{\lambda^2}$.
\end{proposition}

\begin{proof}
Dengan \omterm{thm:g12:integ:ibp}{pengintegralan parsial} pada $\intcc{0}{A}$:
\[
\int_0^A t\,\lambda\eu^{-\lambda t}\dd t
= \bigl[-t\,\eu^{-\lambda t}\bigr]_0^A + \int_0^A \eu^{-\lambda t}\dd t
= -A\eu^{-\lambda A} + \frac{1 - \eu^{-\lambda A}}{\lambda}
\xrightarrow[A\to+\infty]{} \frac1\lambda,
\]
dengan memakai $A\eu^{-\lambda A} \to 0$ (\cref{thm:g12:exp:growth}). \omterm{thm:g12:integ:ibp}{Pengintegralan
parsial} yang kedua memberi $\E(X^2) = \frac{2}{\lambda^2}$, sehingga
$\V(X) = \frac{2}{\lambda^2} - \frac{1}{\lambda^2}$.
\end{proof}

\begin{theorem}[Ketiadaan ingatan]\label{thm:g12:contdist:memoryless}
\index{ketiadaan ingatan}
Jika $X \sim \mathcal E(\lambda)$, maka untuk semua $s, t \geq 0$:
\[
\pcond{X > s}{X > s + t} = \P(X > t) .
\]
\omterm{def:g12:contdist:expo}{Distribusi eksponen} adalah hukum masa hidup \emph{tanpa penuaan}
(inti radioaktif, bukan bola lampu).
\end{theorem}

\begin{proof}
\[
\pcond{X > s}{X > s+t}
= \frac{\P(X > s + t)}{\P(X > s)}
= \frac{\eu^{-\lambda(s+t)}}{\eu^{-\lambda s}} = \eu^{-\lambda t}
= \P(X > t). \qedhere
\]
\end{proof}

\begin{omfigure}
\begin{tikzpicture}
\begin{axis}[omaxis,
  width=8.8cm, height=5cm,
  xmin=-0.4, xmax=4.4, ymin=0, ymax=1.15,
  xtick={1.3}, xticklabels={$t$}, ytick={1}, yticklabels={$\lambda$}]
  \addplot[name path=dens, omDef, thick, domain=0:4.2, samples=80]
    {exp(-x)};
  \path[name path=xaxis] (0,0) -- (4.2,0);
  \addplot[omThm!20] fill between[of=dens and xaxis,
    soft clip={domain=1.3:4.2}];
  \draw[black, dashed, thin] (axis cs:1.3,0) -- (axis cs:1.3,0.2725);
  \node[omThm, font=\small, anchor=west] at (axis cs:2.1,0.35)
    {$\P(X > t) = \eu^{-\lambda t}$};
\end{axis}
\end{tikzpicture}

{\small \omterm{def:g12:contdist:density}{Kerapatan} eksponennya $\lambda\eu^{-\lambda x}$ (di sini
$\lambda = 1$): luas ekornya di seberang $t$ (merah) adalah $\eu^{-\lambda t}$,
dan ketiadaan ingatannya menyatakan bahwa tiap ekornya tampak seperti seluruh \omterm{def:g12:randvar:rv}{distribusinya}
yang diskala ulang.}
\end{omfigure}

\section{Distribusi normal}

\begin{definition}[Distribusi normal]\label{def:g12:contdist:normal}
Peubah $X$ mengikuti \emph{distribusi normal baku}\index{distribusi normal}
$\mathcal N(0, 1)$ jika \omterm{def:g12:contdist:density}{kerapatannya} pada $\R$ adalah
\[
\varphi(t) = \frac{1}{\sqrt{2\pi}}\,\eu^{-t^2/2}
\]
(itulah kurva loncengnya). Secara lebih umum,
$X \sim \mathcal N(\mu, \sigma^2)$ jika $\dfrac{X - \mu}{\sigma} \sim
\mathcal N(0,1)$; lalu $\E(X) = \mu$ dan $\V(X) = \sigma^2$.
\end{definition}

\begin{omfigure}
\begin{tikzpicture}
\begin{axis}[
  width=11cm, height=4.6cm,
  axis lines=middle, xmin=-4, xmax=4, ymin=0, ymax=0.45,
  xtick={-3,-2,-1,0,1,2,3},
  xticklabels={$\mu{-}3\sigma$,$\mu{-}2\sigma$,$\mu{-}\sigma$,$\mu$,
               $\mu{+}\sigma$,$\mu{+}2\sigma$,$\mu{+}3\sigma$},
  ytick=\empty, samples=120, domain=-4:4]
  \addplot[name path=gauss, omDef, thick] {exp(-x^2/2)/sqrt(2*pi)};
  \path[name path=xaxis] (-4,0) -- (4,0);
  \addplot[omDef!20] fill between[of=gauss and xaxis,
           soft clip={domain=-2:2}];
  \node[omDef] at (axis cs:0,0.17) {$\approx 95.4\%$};
\end{axis}
\end{tikzpicture}

{\small \omterm{def:g12:contdist:density}{Kerapatan} normalnya: kira-kira $68.3\%$ massanya terletak dalam jarak
$\sigma$ dari rata-ratanya, $95.4\%$ dalam $2\sigma$, dan $99.7\%$ dalam
$3\sigma$.}
\end{omfigure}

\begin{remark}
Faktor $\frac{1}{\sqrt{2\pi}}$ membuat luas totalnya $1$ --- sebuah perhitungan
yang masyhur (\emph{\omterm{def:g12:integ:area}{integral} Gauss}) yang dikerjakan di universitas. Tak ada
rumus dasar bagi $\int \varphi$: peluang normalnya dibaca dari
tabel atau kalkulator.
\end{remark}

\begin{theorem}[De Moivre--Laplace, diterima saja]\label{thm:g12:contdist:tcl}
\index{teorema De Moivre--Laplace}
Misalkan $X_n \sim \mathcal B(n, p)$ dan
$Z_n = \dfrac{X_n - np}{\sqrt{np(1-p)}}$ (yaitu binomial yang dibakukan). Maka
untuk semua $a \leq b$,
\[
\P(a \leq Z_n \leq b) \xrightarrow[n\to+\infty]{}
\int_a^b \varphi(t)\,\dd t .
\]
\emph{Hasil ini diterima saja pada tingkat ini.} Hasil itu menjelaskan kemunculan
kurva lonceng di mana-mana: binomial dengan $n$ yang besar \omterm{def:g10:numbers:approx}{hampiran} normal
--- dan (menurut teorema limit pusat, di universitas) begitu pula setiap jumlah banyak
pengaruh kecil yang \omterm{def:g12:sums:indep}{saling bebas}.
\end{theorem}

\begin{method}[Selang ayunan pada $95\%$]\label{met:g12:contdist:fluctuation}
\index{selang kepercayaan}
Karena peubah normal jatuh dalam jarak $1.96$ \omterm{def:g11:stat:variance}{simpangan baku} dari rata-ratanya
dengan peluang $0.95$, maka untuk $n$ yang besar frekuensi
$F_n = \frac{X_n}{n}$ dari contoh $\mathcal B(n,p)$ memenuhi
\[
\P\left(p - 1.96\sqrt{\tfrac{p(1-p)}{n}} \leq F_n \leq
p + 1.96\sqrt{\tfrac{p(1-p)}{n}}\right) \approx 0.95 .
\]
\begin{itemize}
  \item \emph{Selang ayunan} (untuk menguji): jika suatu $p$ yang diklaim menempatkan
        frekuensi yang teramati di luar selang ini, tolaklah klaimnya pada
        taraf $5\%$.
  \item \emph{\omterm{met:g12:contdist:fluctuation}{Selang kepercayaan}} (untuk menaksir): selang yang disederhanakan,
        $\left[F_n - \frac{1}{\sqrt n},\ F_n + \frac{1}{\sqrt n}\right]$,
        memuat $p$ dengan peluang sekurang-kurangnya $0.95$
        (dengan memakai $1.96\sqrt{p(1-p)} \leq 1$, lihat \cref{exo:g12:sums:8}).
\end{itemize}
\end{method}

\begin{example}\label{ex:g12:contdist:poll}
Sebuah jajak pendapat atas $n = 1000$ pemilih memberi seorang calon $52\%$. \omterm{met:g12:contdist:fluctuation}{Selang
kepercayaannya} $52\% \pm \frac{1}{\sqrt{1000}} \approx 52\% \pm 3.2\%$ masih
memuat $50\%$: jadi jajak pendapatnya sendirian tidak membuktikan bahwa calonnya
unggul.
\end{example}

\section{Latihan}

\begin{exercise}[$\star$]\label{exo:g12:contdist:1}
Sebuah bus lewat tiap $15$ menit; seorang pelancong tiba pada saat yang acak
seragam. Misalkan $X \sim \mathcal U\intcc{0}{15}$ waktu tunggunya. Hitunglah
$\P(X \leq 5)$, $\P(4 \leq X \leq 10)$, $\E(X)$ dan $\sigma(X)$.
\end{exercise}

\begin{exercise}[$\star$]\label{exo:g12:contdist:2}
Periksalah bahwa $f(t) = \frac{3}{8}t^2$ adalah \omterm{def:g12:contdist:density}{kerapatan} pada $\intcc{0}{2}$, lalu
hitunglah $\E(X)$ dan $\P(X \geq 1)$ bagi peubah $X$ yang berkerapatan ini.
\end{exercise}

\begin{exercise}[$\star$]\label{exo:g12:contdist:3}
Masa hidup (dalam tahun) sebuah komponen elektronik mengikuti
$\mathcal E(0.2)$.
\begin{enumerate}
  \item Hitunglah harapan masa hidupnya dan $\P(X > 5)$.
  \item Komponen itu sudah bekerja $3$ tahun. Berapa
        peluang komponen itu bekerja sekurang-kurangnya $5$ tahun lagi?
\end{enumerate}
\end{exercise}

\begin{exercise}[$\star\star$]\label{exo:g12:contdist:4}
Waktu paruh sebuah atom radioaktif yang masa hidupnya mengikuti
$\mathcal E(\lambda)$ adalah \omterm{def:g11:stat:median}{mediannya} $m$: $\P(X > m) = \frac12$. Nyatakan
$m$ sebagai \omterm{def:g11:func:function}{fungsi} $\lambda$ lalu bandingkan dengan harapannya. Mana yang
lebih besar, dan mengapa itu masuk akal bagi \omterm{def:g12:randvar:rv}{distribusi} yang miring?
\end{exercise}

\begin{exercise}[$\star\star$]\label{exo:g12:contdist:5}
Tinggi badan pada suatu populasi mengikuti $\mathcal N(175,\ 7^2)$ (dalam cm). Dengan
kaidah $68$--$95$--$99.7$, taksirlah bagian populasi yang bertinggi
antara $168$ dan $182$ cm, di atas $189$ cm, dan di bawah $154$ cm.
\end{exercise}

\begin{exercise}[$\star\star$]\label{exo:g12:contdist:6}
Sebuah mesin mengisi kantong berlabel $500$ g; massa yang terisi mengikuti
$\mathcal N(\mu,\ 4^2)$. Peraturannya menuntut agar paling banyak $2.5\%$ kantongnya
berbobot kurang dari $500$ g. Dengan kaidah $95\%$, setelan $\mu$ terkecil apa
yang memenuhinya?
\end{exercise}

\begin{exercise}[$\star\star$]\label{exo:g12:contdist:7}
Sebuah dadu yang diklaim setimbang dilempar $1200$ kali dan menunjukkan angka enam $260$
kali (seperti pada \cref{exo:g12:sums:7}).
\begin{enumerate}
  \item Hitunglah selang ayunan $95\%$ bagi frekuensi angka enam
        pada dadu setimbang sepanjang $1200$ lemparan.
  \item Apakah frekuensi yang teramati berada di dalamnya? Bandingkan kekuatan
        kesimpulan ini dengan telaah Bienaymé--Chebyshev.
\end{enumerate}
\end{exercise}

\begin{exercise}[$\star\star\star$]\label{exo:g12:contdist:8}
Menjelang sebuah pemilu, jajak pendapat atas $n$ orang akan menaksir perolehan $p$ seorang
calon lewat frekuensi yang teramati $F_n$.
\begin{enumerate}
  \item Dengan \omterm{met:g12:contdist:fluctuation}{selang kepercayaan} pada
        \cref{met:g12:contdist:fluctuation}, ukuran contoh berapa yang menjamin
        marjin $\pm 2$ angka?
  \item Kedua calonnya terpisah $1$ angka pada kenyataannya. Jelaskan mengapa
        tak ada jajak pendapat yang realistis dapat menetapkan pemenangnya dengan andal, sebaik
        apa pun jajak itu dijalankan.
\end{enumerate}
\end{exercise}

\begin{exercise}[$\star\star\star$]\label{exo:g12:contdist:9}
Misalkan $X$ berkerapatan $f$ pada $\intcc ab$ dan misalkan
$Y = \alpha X + \beta$ dengan $\alpha > 0$.
\begin{enumerate}
  \item Tunjukkan bahwa $\P(c \leq Y \leq d)
        = \P\left(\frac{c - \beta}{\alpha} \leq X \leq
        \frac{d-\beta}{\alpha}\right)$, lalu simpulkan bahwa $Y$ berkerapatan
        $g(y) = \frac{1}{\alpha} f\left(\frac{y - \beta}{\alpha}\right)$.
  \item Simpulkan bahwa jika $Z \sim \mathcal N(0,1)$, maka
        $\mu + \sigma Z$ berkerapatan
        $\frac{1}{\sigma\sqrt{2\pi}}\,
        \eu^{-(y - \mu)^2/(2\sigma^2)}$ --- yaitu \omterm{def:g12:contdist:density}{kerapatan} normal yang umum.
\end{enumerate}
\end{exercise}

\section{Soal: Menunggu bus, mengukur kerumunan}

\begin{problem}[{Soal akhir pekan --- tiga kerapatan menjalankan
dunia: ketidaktahuan total, penungguan tanpa ingatan, dan lonceng dari banyak
sebab kecil; ditambah paradoks yang membuat busmu selalu terlambat}]
\label{pb:g12:contdist:1}
Mengapa busmu selalu terasa lebih lama daripada yang dijanjikan jadwalnya?
Mengapa kelas yang ditempati murid rata-rata lebih besar
daripada kelas rata-rata? Mengapa temanmu punya lebih banyak teman
daripada kamu? Satu keping matematika yang licik --- pengambilan contoh
yang berbias panjang --- menjawab ketiganya, dan keping itu hidup di bab ini,
di antara \omterm{def:g12:contdist:uniform}{distribusi seragam}, eksponen dan normal
(\cref{def:g12:contdist:uniform}, \cref{def:g12:contdist:expo},
\cref{def:g12:contdist:normal}). Soal penutup jilid
ini mempekerjakan ketiga \omterm{def:g12:contdist:density}{kerapatannya} lalu berakhir di tempat yang ditunjuk seluruh
serinya: kurva lonceng.

\medskip
\noindent\textbf{Bagian I --- Tiga tokoh.}
\begin{enumerate}
  \item Sebuah bus tiba pada saat yang acak seragam dalam
        $30$ menit berikutnya. Hitunglah $\P(\text{tunggu} \leq 10)$, yaitu
        harapan waktu tunggunya, dan $\sigma$
        (\cref{prop:g12:contdist:uniformmoments}).
  \item Periksalah bahwa $f(x) = 2x$ pada $\intcc{0}{1}$ adalah \omterm{def:g12:contdist:density}{kerapatan},
        lalu hitunglah $\E(X)$ dan $\P\left(X > \frac12\right)$.
  \item Panggilan telepon tiba di sebuah saluran bantuan sebagai arus tanpa ingatan:
        waktu tunggu panggilan berikutnya eksponen dengan rata-rata
        $10$ menit. Hitunglah $\P(T > 15)$ dan \omterm{def:g11:stat:median}{median} waktu tunggunya
        --- mengapa \omterm{def:g11:stat:median}{mediannya} \emph{lebih kecil} daripada rata-ratanya?
  \item Kamu sudah menunggu $20$ menit. Berapa
        $\P(T > 35 \mid T > 20)$
        (\cref{thm:g12:contdist:memoryless})? Tafsirkan dalam satu
        kalimat.
  \item Cocokkanlah tiap penungguan dengan model yang tepat --- seragam,
        eksponen, atau bukan keduanya: (a) bunyi klik pencacah
        Geiger; (b) kereta bawah tanah tiap $8$ menit dan
        kedatanganmu tak selaras dengannya; (c) panggilan berikutnya di
        saluran bantuan; (d) putusnya bola lampu yang
        \emph{aus}. Berilah alasan bagi (d) dengan teorema
        ketiadaan ingatannya.
\end{enumerate}

\medskip
\noindent\textbf{Bagian II --- Loncengnya beraksi.}
(Pakailah tetenger normalnya: $68\,\%$ dalam $\sigma$, $95\,\%$
dalam $2\sigma$, $99.7\,\%$ dalam $3\sigma$.)
\begin{enumerate}[resume]
  \item Tinggi badan orang dewasa kira-kira mengikuti
        $\mathcal N(172,\ 8)$ (cm). Berapa bagian yang terletak di
        $\intcc{164}{180}$, $\intcc{156}{188}$,
        $\intcc{148}{196}$?
  \item Taksirlah bagian yang lebih tinggi daripada $190$ cm (\omterm{pb:g11:stat:1}{skor-z}
        dahulu; lalu tabel atau kalkulator).
  \item Skor IQ ditera pada $\mathcal N(100, 15)$.
        Berapa bagian yang bernilai di atas $130$ --- dan kira-kira satu
        orang dari berapa?
  \item Sebuah mesin memotong baut sepanjang
        $\mathcal N(50,\ 0.1)$ mm sedangkan syaratnya
        $\intcc{49.8}{50.2}$. Berapa bagian yang ditolak? Industri
        merayakan proses ``enam sigma'', yang
        syaratnya duduk pada $\pm 6\sigma$: apa yang dibeli hal itu,
        dan mengapa pabrik mengejarnya?
  \item De Moivre--Laplace (\cref{thm:g12:contdist:tcl}): untuk
        $100$ lemparan uang logam yang setimbang, hampirilah
        $\P(45 \leq \text{gambar} \leq 55)$ (dengan kesinambungan:
        $\pm 0.5$; $\sigma = 5$). Mesin kayu yang mana pada
        \cref{pb:g11:binom:1} yang dihaluskan teorema ini menjadi
        sebuah kurva?
  \item Dari \cref{exo:g12:contdist:8}: marjin $\pm 2$ angka
        memerlukan kira-kira $n = 2\,400$. Hitunglah contoh yang diperlukan untuk
        $\pm 0.5$ angka --- lalu simpulkan, dalam satu kalimat, mengapa
        perlombaan yang terpisah satu angka tak dapat ditetapkan secara jujur
        oleh jajak pendapat mana pun.
\end{enumerate}

\medskip
\noindent\textbf{Bagian III --- Paradoks busnya.}
\begin{enumerate}[resume]
  \item Bus tiba sebagai arus tanpa ingatan dengan jarak rata-rata $10$
        menit. Kamu sampai di haltenya pada saat yang sembarang.
        Naluri mengatakan harapan waktu tunggunya $5$ menit
        (yaitu separuh jaraknya). Apa yang justru dikatakan teorema ketiadaan ingatannya?
  \item Bedahlah naluri itu dengan jadwal mainan: jaraknya
        berselang-seling $5$ dan $15$ menit (jarak rata-ratanya $10$).
        Hitunglah peluang bahwa kedatangan yang acak seragam
        mendarat pada jarak yang panjang, lalu harapan waktu tunggu yang sebenarnya ---
        lalu sebutkan biang keladinya: kedatangannya mencuplik jaraknya dengan
        peluang yang sebanding dengan \emph{panjangnya}.
  \item Pada arus tanpa ingatan itu, jarak tempatmu mendarat berharapan
        panjang $20$ menit --- dua kali jarak yang lazim
        (sebab waktu mundur ke bus sebelumnya dan waktu maju ke bus
        berikutnya \emph{keduanya} eksponen dengan rata-rata $10$).
        Damaikanlah ini dengan pertanyaan 12 lalu nyatakan
        \emph{paradoks pemeriksaan} itu dalam satu kalimat.
  \item Sengatan yang sama di darat: sebuah universitas menjalankan sembilan kelas berisi
        $10$ murid dan satu kelas berisi $110$. Hitunglah ukuran kelas
        rata-ratanya, lalu ukuran kelas yang dialami
        \emph{murid rata-rata}. Bilangan mana yang akan dikutip brosurnya,
        dan mana yang menjadi kebenaran yang dijalani muridnya?
  \item Paradoks pertemanan --- ``temanmu, secara
        rata-rata, punya lebih banyak teman daripada kamu'' --- adalah
        matematika yang sama. Katakan dalam satu kalimat apa yang berperan sebagai
        jarak yang panjang itu.
\end{enumerate}

\medskip
\noindent\textbf{Bagian IV --- Menutup jilidnya.}
\begin{enumerate}[resume]
  \item Simulasi, jembatan para praktisi: jika $U$
        seragam pada $\intoo{0}{1}$, tunjukkan bahwa
        $T = -10\ln U$ eksponen dengan rata-rata $10$ (hitunglah
        $\P(T > t)$). Setiap simulasi penungguan acak di
        industri berjalan di atas muslihat satu baris ini.
  \item Resep lama para pemrogram permainan membangun lonceng: jumlahkan
        dua belas peubah seragam yang \omterm{def:g12:sums:indep}{saling bebas} pada $\intcc{0}{1}$ lalu
        kurangkan $6$. Berikan rata-rata dan \omterm{def:g12:randvar:exp}{ragam} jumlahnya,
        lalu jelaskan --- dengan papan Galton di tangan --- mengapa
        hasilnya nyaris normal.
  \item Pasar jatuh lebih keras daripada yang diizinkan loncengnya: keruntuhan tahun
        1987 pernah disebut ``\omterm{def:g11:prob:model}{kejadian} $25\sigma$'', yang
        di bawah kenormalan berpeluang sekitar $10^{-137}$.
        Apakah kesimpulan yang benar --- tentang dunianya, atau
        tentang modelnya? (Satu kalimat, katekismus sang ahli statistika.)
  \item Penutup, dan salam perpisahan bagi jilidnya: ketiga
        tokohnya masing-masing dalam satu baris --- seragam (ketidaktahuan
        dalam batas), eksponen (bahaya tanpa ingatan),
        normal (jumlah banyak sebab kecil); lalu paradoks pemeriksaan
        sebagai sengatan babnya; serta busur
        seluruh serinya --- dari mencacah kardus di kelas 6 sampai
        ke kurva lonceng yang mengukur setiap kerumunan --- dengan
        jilid universitas yang menunggu di tempat limit,
        \omterm{def:g12:integ:area}{integral} dan teorema limit pusat mengambil alih.

\end{enumerate}
\end{problem}
